% BHU PYQ -- Transcribed & Corrected
% Paper: ENGAE-41: Academic Writing
% Exam: B.A. / B.Sc. / B.Com. (Hons.) Semester IV (End Term) Examination, 2025-26 (NEP)
% Category: Ability Enhancement Course (AEC)

\documentclass[11pt,a4paper]{article}
\usepackage[a4paper,top=2.5cm,bottom=2.5cm,left=2.8cm,right=2.8cm,headheight=14pt]{geometry}
\usepackage{amsmath,amssymb}
\usepackage[T1]{fontenc}
\usepackage[utf8]{inputenc}
\usepackage{lmodern,microtype,fancyhdr,enumitem,hyperref,lastpage,parskip}

\hypersetup{
    colorlinks=true,
    urlcolor=black,
    linkcolor=black,
    pdftitle={ENGAE-41: Academic Writing -- BHU Sem IV 2025-26 (NEP)}
}

\pagestyle{fancy}
\fancyhf{}
\renewcommand{\headrulewidth}{0.4pt}
\renewcommand{\footrulewidth}{0.4pt}
\fancyhead[L]{\small   \textit{Banaras Hindu University}}
\fancyhead[R]{\small   \textit{ENGAE-41: Academic Writing (AEC)}}
\fancyfoot[C]{\small Page  \thepage\ of \pageref{LastPage}}


\newlist{parts}{enumerate}{2}
\setlist[parts,1]{label=(\Alph*),leftmargin=*,itemsep=3pt,topsep=2pt}

\newcommand{\pts}[1]{\hfill   \textit{[#1]}}


\begin{document}


\begin{center}
  {\large\bfseries BANARAS HINDU UNIVERSITY}\\[2pt]
  {\normalsize B.A. / B.Sc. / B.Com. (Hons.) --- Semester IV Examination, 2025-26 (NEP)}\\[6pt]
  \rule{0.8\linewidth}{0.5pt}\\[6pt]
  {\large\bfseries Ability Enhancement Course (AEC)}\\[4pt]
  {\large\bfseries Paper Code: ENGAE-41 : Academic Writing}\\[4pt]
  {\normalsize Time: 2 Hours\hfill Maximum Marks: 70}
\end{center}

\rule{\linewidth}{0.4pt}

\smallskip



\noindent   \textit{Instructions: Answer all questions. Each question carries equal marks (1 mark each).}

\bigskip




\noindent   \textbf{Question 1.} What is the primary purpose of Academic Writing?\pts{1}

\begin{parts}
  \item Entertainment and storytelling
  \item Communication of ideas and research
  \item Personal expression only
  \item Social interaction
\end{parts}

\medskip




\noindent   \textbf{Question 2.} Genre-based approaches help students to:\pts{1}

\begin{parts}
  \item Ignore structure
  \item Reduce clarity
  \item Avoid analysis
  \item Understand writing conventions
\end{parts}

\medskip




\noindent   \textbf{Question 3.} Examples of academic genres include:\pts{1}

\begin{parts}
  \item Case studies of different domains
  \item Songs only
  \item Advertisements only
  \item Stories only
\end{parts}

\medskip




\noindent   \textbf{Question 4.} Which is a stage in Process writing?\pts{1}

\begin{parts}
  \item Drafting and revising
  \item Final grading only
  \item Memorizing facts
  \item Copying content
\end{parts}

\medskip




\noindent   \textbf{Question 5.} Online platforms in academic writing encourage:\pts{1}

\begin{parts}
  \item Isolation
  \item Memorization
  \item Collaboration
  \item Avoidance of feedback
\end{parts}

\medskip




\noindent   \textbf{Question 6.} Academic writing in higher education primarily supports:\pts{1}

\begin{parts}
  \item Entertainment
  \item Learning and assessment
  \item Social media use
  \item Personal Journaling
\end{parts}

\medskip




\noindent   \textbf{Question 7.} Integrating approaches means:\pts{1}

\begin{parts}
  \item Using one method only
  \item Combining multiple writing methods
  \item Ignoring feedback
  \item Memorizing facts
\end{parts}

\medskip




\noindent   \textbf{Question 8.} Accuracy in academic writing means:\pts{1}

\begin{parts}
  \item Summarised information
  \item Information is detailed
  \item Information is popular
  \item Information is correct and verifiable
\end{parts}

\medskip




\noindent   \textbf{Question 9.} Overuse of multimedia in academic writing can:\pts{1}

\begin{parts}
  \item Improve clarity always
  \item Distract readers
  \item Replace writing
  \item Increase accuracy
\end{parts}

\medskip




\noindent   \textbf{Question 10.} Multimedia elements in academic documents should be:\pts{1}

\begin{parts}
  \item Irrelevant
  \item Overused
  \item Relevant and purposeful
  \item Flashy and animated
\end{parts}

\medskip




\noindent   \textbf{Question 11.} Identify the correct statement regarding Chicago style citation:\pts{1}

\begin{parts}
  \item MLA uses year in in-text citation
  \item APA uses page number always
  \item Chicago uses notes system
  \item MLA uses footnotes primarily
\end{parts}

\medskip




\noindent   \textbf{Question 12.} Which is the correct Chicago footnote format?\pts{1}

\begin{parts}
  \item Sharma 2021
  \item (Sharma 21)
  \item Sharma,   \textit{Book Title} (2021), 21.
  \item Sharma:21
\end{parts}

\medskip




\noindent   \textbf{Question 13.} Which of the following is a correct MLA in-text citation?\pts{1}

\begin{parts}
  \item (Sharma, 2021)
  \item (Sharma 21)
  \item (Sharma: 2021)
  \item (Sharma p. 21)
\end{parts}

\medskip




\noindent   \textbf{Question 14.} MLA format typically requires which font style and spacing?\pts{1}

\begin{parts}
  \item Times New Roman, 12pt, double-spaced
  \item Arial, 10pt, single-spaced
  \item Courier New, 14pt, double-spaced
  \item Calibri, 11pt, 1.5-spaced
\end{parts}

\medskip




\noindent   \textbf{Question 15.} What is the correct APA in-text citation for a direct quote?\pts{1}

\begin{parts}
  \item (Sharma, 2021, p. 21)
  \item (Sharma 21)
  \item (Sharma: 2021: 21)
  \item (Sharma, 21)
\end{parts}

\medskip




\noindent   \textbf{Question 16.} Academic integrity means:\pts{1}

\begin{parts}
  \item Honest and ethical academic practice
  \item Copying with permission
  \item Using AI without acknowledgment
  \item Buying pre-written papers
\end{parts}

\medskip




\noindent   \textbf{Question 17.} Which of the following is a common citation error?\pts{1}

\begin{parts}
  \item Missing author name
  \item Incorrect formatting
  \item Incomplete references
  \item All of the above
\end{parts}

\medskip




\noindent   \textbf{Question 18.} Which of the following helps to avoid plagiarism?\pts{1}

\begin{parts}
  \item Proper citation
  \item Paraphrasing
  \item Direct quoting with attribution
  \item All of the above
\end{parts}

\medskip




\noindent   \textbf{Question 19.} Plagiarism in academic writing can lead to:\pts{1}

\begin{parts}
  \item Academic success
  \item Rewards
  \item Penalties and disciplinary action
  \item Recognition
\end{parts}

\medskip




\noindent   \textbf{Question 20.} Self-plagiarism refers to:\pts{1}

\begin{parts}
  \item Reusing one's own previously published work without citation
  \item Writing an original paper
  \item Editing someone else's work
  \item Translating a text into another language
\end{parts}

\medskip




\noindent   \textbf{Question 21.} Plagiarism is defined as:\pts{1}

\begin{parts}
  \item Using another person's work or ideas without acknowledgement
  \item Writing a literature review
  \item Summarizing public facts
  \item Citing multiple sources
\end{parts}

\medskip




\noindent   \textbf{Question 22.} Bibliography in Chicago Notes-Bibliography style appears:\pts{1}

\begin{parts}
  \item At the end of the document
  \item In the introduction
  \item At the bottom of each page
  \item In the abstract
\end{parts}

\medskip




\noindent   \textbf{Question 23.} Chicago style is commonly used in which discipline?\pts{1}

\begin{parts}
  \item History and Humanities
  \item Physics
  \item Biology
  \item Mathematics
\end{parts}

\medskip




\noindent   \textbf{Question 24.} Which of the following is NOT a standard citation style?\pts{1}

\begin{parts}
  \item APA
  \item MLA
  \item Chicago
  \item Society of Sciences
\end{parts}

\medskip




\noindent   \textbf{Question 25.} The primary emphasis of a process-oriented approach to writing is on:\pts{1}

\begin{parts}
  \item Final exam scores
  \item Attendance
  \item Writing stages (drafting, revising, editing)
  \item Printing quality
\end{parts}

\medskip




\noindent   \textbf{Question 26.} Toolkit approach in academic writing refers to:\pts{1}

\begin{parts}
  \item Fixed single method
  \item Combining flexible writing strategies and digital tools
  \item Rote memorization
  \item Using physical tools only
\end{parts}

\medskip




\noindent   \textbf{Question 27.} The main purpose of Editing during the writing process is:\pts{1}

\begin{parts}
  \item Idea generation
  \item Checking grammar, spelling, mechanics, and correctness
  \item Brainstorming
  \item Planning paper structure
\end{parts}

\medskip




\noindent   \textbf{Question 28.} Methodology section in a research paper refers to:\pts{1}

\begin{parts}
  \item Primary sources list
  \item Final results
  \item Research methods and procedures used
  \item Scope statement
\end{parts}

\medskip




\noindent   \textbf{Question 29.} Select the correct standard order of sections in a research paper:\pts{1}

\begin{parts}
  \item Abstract -- Introduction -- Methodology -- Results -- Discussion -- Conclusion
  \item Methodology -- Abstract -- Conclusion -- Introduction
  \item Conclusion -- Methodology -- Abstract -- Introduction
  \item Introduction -- Conclusion -- Methodology -- Abstract
\end{parts}

\medskip




\noindent   \textbf{Question 30.} Primary research involves:\pts{1}

\begin{parts}
  \item Direct data collection and original investigation
  \item Reading secondary summaries
  \item Copying textbook chapters
  \item Reviewing Wikipedia entries
\end{parts}

\medskip




\noindent   \textbf{Question 31.} Primary sources do NOT include:\pts{1}

\begin{parts}
  \item Articles summarizing other political opinion pieces
  \item Statistical raw data
  \item Original archival letters
  \item Experimental observation logs
\end{parts}

\medskip




\noindent   \textbf{Question 32.} In MLA style, entries in the Works Cited page are ordered by:\pts{1}

\begin{parts}
  \item Alphabetical order of author's last name
  \item Order of appearance in text
  \item Publication date order
  \item Page length of sources
\end{parts}

\medskip




\noindent   \textbf{Question 33.} While preparing a slide presentation (PPT), it is best practice to:\pts{1}

\begin{parts}
  \item Fill the slide with dense paragraphs
  \item Use very flashy neon colors
  \item Use bullet points and clear visuals
  \item Use images with no text
\end{parts}

\medskip




\noindent   \textbf{Question 34.} Which tool is commonly used for digital academic presentations?\pts{1}

\begin{parts}
  \item Microsoft PowerPoint / Google Slides
  \item Microsoft Teams chat
  \item Notepad
  \item Paint
\end{parts}

\medskip




\noindent   \textbf{Question 35.} Multimedia integration is easiest and most interactive in:\pts{1}

\begin{parts}
  \item Print Journals
  \item Digital academic writing and e-publications
  \item Field notebook entries
  \item Handwritten reports
\end{parts}

\medskip




\noindent   \textbf{Question 36.} The main job of a topic sentence in a paragraph is to:\pts{1}

\begin{parts}
  \item State the main idea of the paragraph
  \item Conclude the entire paper
  \item List all references
  \item Provide a footnote
\end{parts}

\medskip




\noindent   \textbf{Question 37.} A major advantage of digital writing tools is:\pts{1}

\begin{parts}
  \item Easy formatting, editing, and real-time revision
  \item Slower editing and revision
  \item Limited access to documents
  \item No collaboration features
\end{parts}

\medskip




\noindent   \textbf{Question 38.} Which word best demonstrates cautious hedging in academic writing?\pts{1}

\begin{parts}
  \item Demonstrates
  \item Suggests
  \item Proves
  \item Confirms
\end{parts}

\medskip




\noindent   \textbf{Question 39.} Multimedia in digital academic writing includes:\pts{1}

\begin{parts}
  \item Text only
  \item Numbers only
  \item Images, audio, and embedded video
  \item Footnotes only
\end{parts}

\medskip




\noindent   \textbf{Question 40.} A product-oriented approach to writing focuses on:\pts{1}

\begin{parts}
  \item Writing process stages
  \item Brainstorming sessions
  \item Final written output
  \item Peer feedback loops
\end{parts}

\medskip




\noindent   \textbf{Question 41.} A process-oriented approach emphasizes:\pts{1}

\begin{parts}
  \item Steps like drafting, feedback, and revising
  \item Final grade only
  \item Memorization of facts
  \item Copying model texts
\end{parts}

\medskip




\noindent   \textbf{Question 42.} Genre-based writing instruction focuses on:\pts{1}

\begin{parts}
  \item Understanding social context and text conventions
  \item Grammar drills only
  \item Spelling tests
  \item Speed typing
\end{parts}

\medskip




\noindent   \textbf{Question 43.} Learners using a toolkit approach:\pts{1}

\begin{parts}
  \item Follow one rigid method only
  \item Combine various flexible writing strategies
  \item Plagiarise content
  \item Ignore peer feedback
\end{parts}

\medskip




\noindent   \textbf{Question 44.} Academic writing helps students to:\pts{1}

\begin{parts}
  \item Memorize textbook content
  \item Avoid critical analysis
  \item Develop critical thinking and analytical reasoning
  \item Copy information quickly
\end{parts}

\medskip




\noindent   \textbf{Question 45.} Integrating writing approaches means:\pts{1}

\begin{parts}
  \item Using one single method exclusively
  \item Combining multiple writing methods effectively
  \item Rejecting instructor feedback
  \item Avoiding revision
\end{parts}

\medskip




\noindent   \textbf{Question 46.} Peer review in academic writing means:\pts{1}

\begin{parts}
  \item Teacher correction only
  \item Self-editing without help
  \item Constructive feedback from classmates or peers
  \item Literature review search
\end{parts}

\medskip




\noindent   \textbf{Question 47.} Which of the following is a primary characteristic of formal academic tone?\pts{1}

\begin{parts}
  \item Use of slang and informal contractions
  \item Objective, precise, and evidence-based language
  \item Heavy use of personal anecdotes
  \item Emotional language
\end{parts}

\medskip




\noindent   \textbf{Question 48.} An Abstract in a research paper serves to:\pts{1}

\begin{parts}
  \item Provide a concise summary of the entire study
  \item List all cited authors in full
  \item Give personal biographical details
  \item Present raw survey data
\end{parts}

\medskip




\noindent   \textbf{Question 49.} In APA style (7th edition), the reference list is ordered:\pts{1}

\begin{parts}
  \item Alphabetically by author's last name
  \item Chronologically by year
  \item By order of mention in text
  \item By title length
\end{parts}

\medskip




\noindent   \textbf{Question 50.} Which section of a paper interprets findings and compares them with previous studies?\pts{1}

\begin{parts}
  \item Introduction
  \item Discussion
  \item Title Page
  \item Abstract
\end{parts}

\medskip




\noindent   \textbf{Question 51.} Multimedia integration is easier in:\pts{1}

\begin{parts}
  \item Print Journal
  \item Digital academic writing and e-publications
  \item Field work notebook
  \item Hand written reports
\end{parts}

\medskip




\noindent   \textbf{Question 52.} The main job of a topic sentence is to:\pts{1}

\begin{parts}
  \item Ask a question for the reader only
  \item Provide a citation
  \item Summarize the whole article
  \item Introduce the key idea of the paragraph
\end{parts}

\medskip




\noindent   \textbf{Question 53.} Which of the following is a digital writing tool?\pts{1}

\begin{parts}
  \item Notebook
  \item Chalkboard
  \item Microsoft Word
  \item Blackboard (physical)
\end{parts}

\medskip




\noindent   \textbf{Question 54.} LaTeX is mainly used for:\pts{1}

\begin{parts}
  \item Graphic design
  \item Academic and technical writing
  \item Video editing
  \item Email writing
\end{parts}

\medskip




\noindent   \textbf{Question 55.} A major advantage of digital writing tools is:\pts{1}

\begin{parts}
  \item Easy formatting, editing, and real-time revision
  \item Slower editing and revision
  \item Limited access to documents
  \item No collaboration features
\end{parts}

\medskip




\noindent   \textbf{Question 56.} Which word best shows caution (hedging) in academic writing?\pts{1}

\begin{parts}
  \item Demonstrates
  \item Suggests
  \item Proves
  \item Confirms
\end{parts}

\medskip




\noindent   \textbf{Question 57.} Multimedia in digital academic writing includes:\pts{1}

\begin{parts}
  \item Text only
  \item Numbers only
  \item Images, audio, and video
  \item Footnotes only
\end{parts}

\medskip




\noindent   \textbf{Question 58.} Which is an example of multimedia integration?\pts{1}

\begin{parts}
  \item Writing only text
  \item Embedding a video in a presentation
  \item Printing notes
  \item Using pen and paper
\end{parts}

\medskip




\noindent   \textbf{Question 59.} A 'literature review' primarily aims to:\pts{1}

\begin{parts}
  \item Summarize and synthesize existing research on a topic
  \item Give your personal opinion
  \item List the works you have read
  \item Provide a paraphrase
\end{parts}

\medskip




\noindent   \textbf{Question 60.} The changing landscape of academic writing is mainly influenced by:\pts{1}

\begin{parts}
  \item Traditional methods only
  \item Oral traditions
  \item Handwritten assignments
  \item Technological advancements
\end{parts}

\medskip




\noindent   \textbf{Question 61.} Digital collaboration in academic writing offers:\pts{1}

\begin{parts}
  \item Real-time co-authoring and editing
  \item Less access to sources
  \item Less feedback
  \item Isolation
\end{parts}

\medskip




\noindent   \textbf{Question 62.} A product-oriented approach focuses on:\pts{1}

\begin{parts}
  \item Writing process
  \item Brainstorming
  \item Final written output
  \item Peer feedback
\end{parts}

\medskip




\noindent   \textbf{Question 63.} A process-oriented approach emphasizes:\pts{1}

\begin{parts}
  \item Steps like drafting, revising, and feedback
  \item Final grade only
  \item Memorization
  \item Copying
\end{parts}

\medskip




\noindent   \textbf{Question 64.} Genre-based writing instruction focuses on:\pts{1}

\begin{parts}
  \item Grammar only
  \item Conceptualizing
  \item Random writing
  \item Specific text types and conventions
\end{parts}

\medskip




\noindent   \textbf{Question 65.} The toolkit approach includes:\pts{1}

\begin{parts}
  \item Multiple tools and strategies
  \item Only one method
  \item No resources
  \item Only lectures
\end{parts}

\medskip




\noindent   \textbf{Question 66.} Learners using a toolkit approach:\pts{1}

\begin{parts}
  \item Follow one method only
  \item Combine various writing strategies
  \item Plagiarise
  \item Ignore feedback
\end{parts}

\medskip




\noindent   \textbf{Question 67.} Academic writing helps students to:\pts{1}

\begin{parts}
  \item Memorize content
  \item Avoid analysis
  \item Develop critical thinking
  \item Copy information
\end{parts}

\medskip




\noindent   \textbf{Question 68.} Integrating approaches means:\pts{1}

\begin{parts}
  \item Using one method only
  \item Combining multiple writing methods
  \item Avoiding strategies
  \item Ignoring structure
\end{parts}

\medskip




\noindent   \textbf{Question 69.} Which of the following is a secondary resource for a researcher in commerce?\pts{1}

\begin{parts}
  \item Articles on Business case studies
  \item A live scientific experiment
  \item Historical documents
  \item Works of Art
\end{parts}

\medskip




\noindent   \textbf{Question 70.} Peer review means:\pts{1}

\begin{parts}
  \item Teacher correction
  \item Self-editing
  \item Feedback from peers and colleagues
  \item Literature review search
\end{parts}

\vfill
\rule{\linewidth}{0.4pt}

\begin{center}\small
\href{https://bhu-exam.vercel.app/}{Click here to visit our website}\\[4pt]
\href{https://me-aryan.vercel.app/}{Click here to visit our contributor}
\end{center}

\end{document}
